%    For use when working on individual chapters
%\includeonly{}

%    Include referenced packages here.
\usepackage[left =.25in, right = .25in, top = 1in, bottom = 1in]{geometry} 
\usepackage{amsmath}
\usepackage{amsthm}
\usepackage{amscd}
\usepackage{amssymb}
\usepackage{setspace}
\usepackage{mathtools}
\usepackage{tikz}  
\usepackage{tikz-cd}
\usepackage{tkz-fct}
\usepackage{amsrefs}

\usepackage{pgfplots}
\pgfplotsset{compat=1.18}

\usepackage{environ}
\usepackage{tikz-cd} 
\usepackage{enumitem}
\usepackage{xcolor}   %May be necessary if you want to color links
%\usepackage{xr}

\usepackage{hyperref}
\hypersetup{
	colorlinks=true, %set true if you want colored links
	linktoc=all,     %set to all if you want both sections and subsections linked
	linkcolor=black,  %choose some color if you want links to stand out
	urlcolor=cyan
}
\usepackage[symbols,nogroupskip,sort=none]{glossaries-extra}

\pgfplotsset{every axis/.append style={
		axis x line=middle,    % put the x axis in the middle
		axis y line=middle,    % put the y axis in the middle
		axis line style={<->,color=black}, % arrows on the axis
		xlabel={$x$},          % default put x on x-axis
		ylabel={$y$},          % default put y on y-axis
}}


\theoremstyle{definition}
\newtheorem{definition}{Definition}[subsection]
\newtheorem{defn}[definition]{Definition}
\newtheorem{note}[definition]{Note}
\newtheorem{ax}[definition]{Axiom}
\newtheorem{thm}[definition]{Theorem}
\newtheorem{lem}[definition]{Lemma}
\newtheorem{prop}[definition]{Proposition}
\newtheorem{cor}[definition]{Corollary}
\newtheorem{conject}[definition]{Conjecture}
\newtheorem{ex}[definition]{Exercise}
\newtheorem{exmp}[definition]{Example}
\newtheorem{soln}[definition]{Solution}
\newtheorem{questn}[definition]{Question}

\setcounter{tocdepth}{3}

% hide proofs
\newif\ifhideproofs
%\hideproofstrue %uncomment to hide proofs
\ifhideproofs
\NewEnviron{hide}{}
\let\proof\hide
\let\endproof\endhide
\fi

% lower-case greek
\newcommand{\al}{\alpha}
\newcommand{\be}{\beta}
\newcommand{\gam}{\gamma}
\newcommand{\del}{\delta}
\newcommand{\ep}{\epsilon}
\newcommand{\ze}{\zeta} 
\renewcommand{\th}{\theta}
\newcommand{\kap}{\kappa} 
\newcommand{\lam}{\lambda}  
\newcommand{\sig}{\sigma} 
\newcommand{\omi}{\omicron}
\newcommand{\up}{\upsilon}
\newcommand{\om}{\omega}

% upper-case greek
\newcommand{\Gam}{\Gamma}
\newcommand{\Del}{\Delta}
\newcommand{\Lam}{\Lambda} 
\newcommand{\Sig}{\Sigma} 
\newcommand{\Om}{\Omega}


% text
\newcommand{\tA}{\text{A}}
\newcommand{\tB}{\text{B}}
\newcommand{\tC}{\text{C}}
\newcommand{\tD}{\text{D}}
\newcommand{\tE}{\text{E}}


% blackboard bold
\newcommand{\C}{\mathbb{C}}
\newcommand{\E}{\mathbb{E}}
\newcommand{\F}{\mathbb{F}}
\renewcommand{\H}{\mathbb{H}}
\newcommand{\K}{\mathbb{K}}
\newcommand{\N}{\mathbb{N}}
\renewcommand{\O}{\mathbb{O}}
\newcommand{\Q}{\mathbb{Q}}
\newcommand{\R}{\mathbb{R}}
\renewcommand{\S}{\mathbb{S}}
\newcommand{\T}{\mathbb{T}}
\newcommand{\V}{\mathbb{V}}
\newcommand{\Z}{\mathbb{Z}}

% math caligraphic
\newcommand{\MA}{\mathcal{A}}
\newcommand{\MB}{\mathcal{B}}
\newcommand{\MC}{\mathcal{C}}
\newcommand{\MD}{\mathcal{D}}
\newcommand{\ME}{\mathcal{E}}
\newcommand{\MF}{\mathcal{F}}
\newcommand{\MG}{\mathcal{G}}
\newcommand{\MH}{\mathcal{H}}
\newcommand{\MI}{\mathcal{I}}
\newcommand{\MJ}{\mathcal{J}}
\newcommand{\MK}{\mathcal{K}}
\newcommand{\ML}{\mathcal{L}}
\newcommand{\MM}{\mathcal{M}}
\newcommand{\MN}{\mathcal{N}}
\newcommand{\MO}{\mathcal{O}}
\newcommand{\MP}{\mathcal{P}}
\newcommand{\MQ}{\mathcal{Q}}
\renewcommand{\MR}{\mathcal{R}}
\newcommand{\MS}{\mathcal{S}}
\newcommand{\MT}{\mathcal{T}}
\newcommand{\MU}{\mathcal{U}}
\newcommand{\MV}{\mathcal{V}}
\newcommand{\MW}{\mathcal{W}}
\newcommand{\MX}{\mathcal{X}}
\newcommand{\MY}{\mathcal{Y}}
\newcommand{\MZ}{\mathcal{Z}}

% mathfrak
\newcommand{\MFX}{\mathfrak{X}}
\newcommand{\MFg}{\mathfrak{g}}
\newcommand{\MFh}{\mathfrak{h}}

% tilde 
\newcommand{\tMA}{\tilde{\MA}}
\newcommand{\tMB}{\tilde{\MB}}
\newcommand{\tU}{\tilde{U}}
\newcommand{\tV}{\tilde{V}}
\newcommand{\tphi}{\tilde{\phi}}
\newcommand{\tpsi}{\tilde{\psi}}
\newcommand{\tF}{\tilde{F}}

\newcommand{\iid}{\stackrel{iid}{\sim}}

\newcommand{\Cn}{\text{Cn}} % cone base shift functor

\newcommand{\cur}{\text{cur}} % curry operator





% label/reference
% internal label/reference
\newcommand{\lex}[1]{\label{ex:#1}}
\newcommand{\rex}[1]{Exercise \ref{ex:#1}}

\newcommand{\ld}[1]{\label{defn:#1}}
\newcommand{\rd}[1]{Definition \ref{defn:#1}}

\newcommand{\lnt}[1]{\label{note:#1}}
\newcommand{\rnt}[1]{Note \ref{note:#1}}

\newcommand{\lax}[1]{\label{ax:#1}}
\newcommand{\rax}[1]{Axiom \ref{ax:#1}}

\newcommand{\lfig}[1]{\label{fig:#1}}
\newcommand{\rfig}[1]{Figure \ref{fig:#1}}

\newcommand{\lsect}[1]{\label{sect:#1}}
\newcommand{\rsect}[1]{Section \ref{sect:#1}}
\newcommand{\rsectp}[1]{(Section \ref{sect:#1})}

% external reference
\newcommand{\extrex}[2]{Exercise \ref{#1-ex:#2}}

\newcommand{\extrd}[2]{Definition \ref{#1-defn:#2}}

\newcommand{\extrax}[2]{Axiom \ref{#1-ax:#2}}

\newcommand{\extrfig}[2]{Figure \ref{#1-fig:#2}}



% external documents (EDIT HERE)
%\externaldocument[analysis-]{"/home/carson/Desktop/Github/Mathematics/Introduction to Analysis/Introduction to Analysis.tex"}











% math operators
\DeclareMathOperator{\supp}{supp} % support
\DeclareMathOperator{\sgn}{sgn} % signum
\DeclareMathOperator{\spn}{span} % span

\DeclareMathOperator{\Obj}{Obj} % objects
\DeclareMathOperator{\dom}{dom} % domain
\DeclareMathOperator{\cod}{cod} % codomain

\DeclareMathOperator*{\ot}{\otimes} % otimes













\newcommand{\Oper}{\text{Op}}
\newcommand{\Ext}{\text{Ext}} % extensions
\newcommand{\Rstr}{\text{Rstr}} % restrictions
\newcommand{\OperExt}{\text{OpExt}}
\newcommand{\unb}{\text{unb}}























\DeclareMathOperator{\Eq}{Eq}

\DeclareMathOperator{\Sym}{Sym}
\DeclareMathOperator{\Alt}{Alt}









% linear algebra
\newcommand{\Dual}{\text{Dual}}
% ######################################################
% ######################################################
% ######################################################


































% ###################_________Sets_________###################
\newcommand{\Cover}{\text{Cover}} % cover

\DeclareMathOperator*{\prodrel}{\Pi^{\Rel}} % product of relations 
\DeclareMathOperator*{\timesrel}{\times^{\Rel}} % binary product of relations  

% Sets - Classes
\newcommand{\Set}{\text{Set}} % class of sets
\newcommand{\Rel}{\text{Rel}} % class of relations

% Sets - Categories
\newcommand{\bSet}{\text{\tbf{Set}}} % category of sets
\newcommand{\bRel}{\text{\tbf{Rel}}} % category of relations
% ######################################################
% ######################################################
% ######################################################
















% ###################_________Logic_________###################

\DeclareMathOperator{\Symb}{\text{Symb}} % set of symbols
\DeclareMathOperator{\RelSymb}{\text{RelSymb}} % set of relation symbols
\DeclareMathOperator{\FunSymb}{\text{FunSymb}} % set of function sympbols
\DeclareMathOperator{\ConsSymb}{\text{ConsSymb}}  % set of constant symbols
\DeclareMathOperator{\ConnSymb}{\text{ConnSymb}}  % set of connective symbols
\DeclareMathOperator{\Quan}{\text{Quan}}  % 
\DeclareMathOperator{\AuxSymb}{\text{AuxSymb}}  % set of auxiliary symbols
\DeclareMathOperator{\Term}{\text{Term}}  % set of terms
\DeclareMathOperator{\ElmTerm}{\text{ElmTerm}}  % elementary terms
\DeclareMathOperator{\AtomForm}{\text{AtomForm}} % atomic formulae

\DeclareMathOperator{\existsop}{\exists} % exists (operator)
% ######################################################
% ######################################################
% ######################################################























% ###################_________Topology_________###################

\DeclareMathOperator{\Homeo}{\text{Homeo}}

\DeclareMathOperator{\cl}{cl} % closure
\DeclareMathOperator{\Int}{Int} % interior
\DeclareMathOperator{\clset}{clSets} % closed supersets
\DeclareMathOperator{\Intset}{IntSets} % open subsets
\DeclareMathOperator*{\Punct}{\text{Punct}} 
\newcommand{\opn}[1]{{#1}^{\text{opn}}} % open adjective
\DeclareMathOperator{\FIP}{\text{FIP}} % finite intersection property

% Metric Spaces 
\newcommand{\Bnd}{\text{Bnd}} % set of bounded functions
\newcommand{\SemiNorm}{\text{SemiNorm}} % set of seminorms
\newcommand{\taulc}{\tau^{\text{lc}}} 


\newcommand{\Lip}{\text{Lip}}  % set of Lipschitz functions
\newcommand{\LipConst}{\text{LipConst}}  % set of Lipschitz constants of a function

\newcommand{\lip}{\text{lip}}  % Lipschitz constant of a function

\newcommand{\diln}{\text{diln}}  % dilation of a function 

\newcommand{\Alx}{\text{Alx}}

% Classes
\DeclareMathOperator*{\Top}{\text{Top}}
\DeclareMathOperator*{\TopEq}{\tbf{TopEq}}
\DeclareMathOperator*{\Frm}{\tbf{Frm}}
\DeclareMathOperator*{\Loc}{\tbf{Loc}}  

% Categories
\DeclareMathOperator*{\bTop}{\text{\tbf{Top}}} % category of topological spaces 
\DeclareMathOperator*{\bTopEq}{\text{\tbf{TopEq}}} % category of topological spaces with equivalence relations
\DeclareMathOperator*{\bFrm}{\text{\tbf{Frm}}} % category of frames
\DeclareMathOperator*{\bLoc}{\text{\tbf{Loc}}} % category of Locales
% ######################################################
% ######################################################
% ######################################################











% ###################_________Measure and Integration_________###################
\newcommand{\sigAlg}{\sig\text{-Alg}}

\newcommand{\Darb}{\text{Darb}} % Darboux integrable functions
\newcommand{\Riem}{\text{Riem}} % Riemann integrable functions
\newcommand{\IntSet}{\text{IntSet}}
\newcommand{\Tag}{\text{Tag}}
\newcommand{\mesh}{\text{mesh}}

\newcommand{\dt}{\, d t}
\newcommand{\ds}{\, d s}
\newcommand{\dx}{\, d x}

\newcommand{\ddel}{\, d \del}
\newcommand{\dmu}{\, d \mu}
\newcommand{\dnu}{\, d \nu}
\newcommand{\dlam}{\, d \lambda}
\newcommand{\dP}{\, d P}
\newcommand{\dQ}{\, d Q}
\newcommand{\dm}{\, d m}
\newcommand{\dsh}{\, d \#}


% Classes
\DeclareMathOperator{\Meas}{\text{Meas}}
\DeclareMathOperator{\TopMeas}{\text{TopMeas}}

\DeclareMathOperator{\Msr}{\text{Msr}}
\DeclareMathOperator{\Msrone}{\text{Msr}_1} 
\DeclareMathOperator{\MsrC}{\text{Msr}_{\C}} 

\DeclareMathOperator{\RadMsr}{\text{RadMsr}}
\DeclareMathOperator{\RadMsrone}{\text{RadMsr}_1}
\DeclareMathOperator{\RadMsrC}{\text{RadMsr}_{\C}}

\DeclareMathOperator{\TopMsr}{\text{TopMsr}}
\DeclareMathOperator{\TopMsrone}{\text{TopMsr}_1} 
\DeclareMathOperator{\TopMsrC}{\text{TopMsr}_{\C}} 

\DeclareMathOperator{\TopRadMsr}{\text{TopRadMsr}}
\DeclareMathOperator{\TopRadMsrone}{\text{TopRadMsr}_1}
\DeclareMathOperator{\TopRadMsrC}{\text{TopRadMsr}_{\C}} 


% Categories
\DeclareMathOperator{\bMeas}{\text{\tbf{Meas}}}
\DeclareMathOperator{\bTopMeas}{\text{\tbf{TopMeas}}}

\DeclareMathOperator{\bMsr}{\text{\tbf{Msr}}} 
\DeclareMathOperator{\bMsrone}{\text{\tbf{Msr}}_1} 
\DeclareMathOperator{\bMsrC}{\text{\tbf{Msr}}_{\C}} 

\DeclareMathOperator{\bRadMsr}{\text{\tbf{RadMsr}}}
\DeclareMathOperator{\bRadMsrone}{\text{\tbf{RadMsr}}_1}
\DeclareMathOperator{\bRadMsrC}{\text{\tbf{RadMsr}}_{\C}}

\DeclareMathOperator{\bTopMsr}{\text{\tbf{TopMsr}}} 
\DeclareMathOperator{\bTopMsrone}{\text{\tbf{TopMsr}}_1} 
\DeclareMathOperator{\bTopMsrC}{\text{\tbf{TopMsr}}_{\C}} 

\DeclareMathOperator{\bTopRadMsr}{\text{\tbf{TopRadMsr}}}
\DeclareMathOperator{\bTopRadMsrone}{\text{\tbf{TopRadMsr}}_1}
\DeclareMathOperator{\bTopRadMsrC}{\text{\tbf{TopRadMsr}}_{\C}} 
% ######################################################
% ######################################################
% ######################################################


































% ###################_________Probability_________###################

\newcommand{\Filtn}{\text{Filtrn}}
% ######################################################
% ######################################################
% ######################################################













































% ###################_________Geometry_________###################

\DeclareMathOperator{\Atlasinf}{\text{Atlasinf}} % collection of smooth atlases on a topological manifold
\DeclareMathOperator{\Diffeo}{\text{Diffeo}}


\newcommand{\germ}{\text{germ}} % germ of functions at at point
\newcommand{\germinf}{\text{germ}^{\infty}} % germ of functions at at point

\newcommand{\Paths}{\text{Paths}} % Paths of a manifold

\newcommand{\LocTriv}{\text{LocTriv}} % local trivializations
\newcommand{\LocTrivz}{\text{LocTriv}^0} % topological local trivializations
\newcommand{\LocTrivinf}{\text{LocTriv}^{\infty}} % smooth local trivializations

\newcommand{\Lie}{\text{Lie}} % Lie algebra of a Lie group


% Geom - Classes
\newcommand{\Manz}{\text{Man}^0} % class of topological manifolds
\newcommand{\Maninf}{\text{Man}^{\infty}}  % class of smooth manifolds with empty boundary
\newcommand{\ManBndinf}{\text{ManBnd}^{\infty}} % class of manifolds with nonempty boundary 

\newcommand{\LieGrp}{\text{LieGrp}} % class of Lie groups


% Geom - Categories
\DeclareMathOperator*{\bManz}{\text{\tbf{Man}}^0} % cat of topological manifolds
\DeclareMathOperator*{\bManinf}{\text{\tbf{Man}}^{\infty}} % cat of smooth manifolds with empty boundary
\DeclareMathOperator*{\bManBndinf}{\text{\tbf{ManBnd}}^{\infty}} % cat of manifolds with nonempty boundary 

\DeclareMathOperator*{\bBunz}{\text{\tbf{Bun}}^0} % cat of topological fiber bundles
\DeclareMathOperator*{\bBuninf}{\text{\tbf{Bun}}^{\infty}} % cat of smooth fiber bundles
\DeclareMathOperator*{\bVecBuninf}{\text{\tbf{VecBun}}^{\infty}} % cat of smooth vector bundles
\DeclareMathOperator*{\bPrinBuninf}{\text{\tbf{PrinBun}}^{\infty}} % cat of smooth principle bundles 

\DeclareMathOperator*{\bLieGrp}{\text{\tbf{LieGrp}}} 



\DeclareMathOperator{\Conn}{\text{Conn}}  % set of connections
% ######################################################
% ######################################################
% ######################################################

















% ###################_________Convexity_________###################

\DeclareMathOperator{\bal}{bal} % balanced subset generated by
\DeclareMathOperator{\cyc}{cyc} % 
\DeclareMathOperator{\cnv}{conv} % convex subset generate by


\DeclareMathOperator{\cnvspn}{\text{convspan}} 
\DeclareMathOperator{\extrm}{extrm} 


\DeclareMathOperator{\epi}{epi} % epigraph 
\DeclareMathOperator{\env}{env} % convex envelope

\newcommand{\Conv}{\text{Conv}} % convex subsets
\newcommand{\Bal}{\text{Bal}} % balanced subsets
\newcommand{\Aff}{\text{Aff}} % affine functions

\newcommand{\Cone}{\text{Cone}} % set of cones
\newcommand{\ConvCone}{\text{ConvCone}} % set of convex cones
\newcommand{\PtConvCone}{\text{PtConvCone}} % set of pointed convex cone
% ######################################################
% ######################################################
% ######################################################
 



\DeclareMathOperator{\codim}{codim} % codimension

\DeclareMathOperator{\Derivinf}{Deriv^{\infty}}

\DeclareMathOperator*{\argmax}{arg\,max}  
\DeclareMathOperator*{\argmin}{arg\,min}
\DeclareMathOperator{\diam}{\text{diam}} % diameter
\DeclareMathOperator{\rnk}{\text{rank}} % rank
\DeclareMathOperator{\tr}{\text{tr}} % trace
\DeclareMathOperator{\prj}{\text{proj}} % projection
\DeclareMathOperator{\nab}{\nabla} % nabla
\DeclareMathOperator{\diag}{\text{diag}} % diagonal
\DeclareMathOperator*{\ind}{\text{ind}} % 

\newcommand{\Part}{\text{Part}} 

\DeclareMathOperator{\Var}{\text{Var}}













































\DeclareMathOperator{\Subcat}{Subcat} % SubCat
\DeclareMathOperator{\SubcatGen}{SubcatGen} % subcategory generated by morphisms
\newcommand{\leqcat}{\leq_\text{Cat}} % Subcat partial order

\DeclareMathOperator{\Subquiv}{Subquiv} % Subquiv
\DeclareMathOperator{\SubquivGen}{\text{SubquivGen}} % subquiver generated by vertices/edges
\newcommand{\leqquiv}{\leq_\text{Quiv}} % Subquiver partial order


\DeclareMathOperator{\EmbSubmaninf}{\text{EmbSubman}^{\infty}} % Subquiv
\DeclareMathOperator{\EmbSubmaninfGen}{\text{EmbSubman}^{\infty}\text{Gen}} % subquiver generated by vertices/edges
\newcommand{\leqembmaninf}{\leq_{\text{Man}^{\infty}}} % Subquiver partial order

\DeclareMathOperator{\lmt}{lim}





























% ###################_________Algebra_________###################


% Groups, Rings, Fields, Modules, Algebras
\newcommand{\inv}{\text{inv}}
\newcommand{\mult}{\text{mult}}
\newcommand{\add}{\text{add}}
\newcommand{\scalmult}{\text{scalar\textunderscore mult}}
%\DeclareMathOperator{\smult}{\text{smult}}

\DeclareMathOperator{\End}{\text{End}} 
\DeclareMathOperator{\rep}{\text{Rep}} 
\DeclareMathOperator*{\Equi}{\text{Equiv}}
\DeclareMathOperator*{\Cong}{\text{Cong}}



% Classes
\newcommand{\Mon}{\text{Mon}} % class of monoids
\newcommand{\Grp}{\text{Grp}} % class of groups
\newcommand{\Semgrp}{\text{Semgrp}} % class of semigroups

\newcommand{\Ring}{\text{Ring}} % class of rings
\newcommand{\Field}{\text{Field}} % class of fields

\newcommand{\Mod}{\text{Mod}} % class of modules

\newcommand{\Vect}{\text{Vect}} % class of vector spaces
\newcommand{\VectR}{\text{Vect}_{\R}} % class of real vector spaces
\newcommand{\VectC}{\text{Vect}_{\C}} % class of complex vector spaces
\newcommand{\VectK}{\text{Vect}_{\K}} % class of K-vector spaces

\newcommand{\Alg}{\text{Alg}} % class of algebras
\newcommand{\AlgR}{\text{Alg}_{\R}} % class of real algebras
\newcommand{\AlgC}{\text{Alg}_{\C}} % class of complex algebras
\newcommand{\AlgK}{\text{Alg}_{\K}} % class of K-algebras



% Categories
\DeclareMathOperator*{\bMon}{\text{\tbf{Mon}}} % cat of monoids
\DeclareMathOperator*{\bGrp}{\text{\tbf{Grp}}} % cat of groups
\DeclareMathOperator*{\bSemgrp}{\text{\tbf{Semgrp}}} % cat of semigroups

\DeclareMathOperator*{\bRing}{\text{\tbf{Ring}}} % cat of rings
\DeclareMathOperator*{\bField}{\text{\tbf{Field}}} % cat of fields

\DeclareMathOperator*{\bMod}{\text{\tbf{Mod}}} % cat of modules

\DeclareMathOperator*{\bVect}{\text{\tbf{Vect}}} % cat of vector spaces
\DeclareMathOperator*{\bVectR}{\text{\tbf{Vect}}_{\R}} % cat of real vector spaces
\DeclareMathOperator*{\bVectC}{\text{\tbf{Vect}}_{\C}} % cat of complex vector spaces
\DeclareMathOperator*{\bVectK}{\text{\tbf{Vect}}_{\K}} % cat of K-vector spaces

\DeclareMathOperator*{\bAlg}{\text{\tbf{Alg}}} % cat of algebras
\DeclareMathOperator*{\bAlgR}{\text{\tbf{Alg}}_{\R}} % cat of real algebras
\DeclareMathOperator*{\bAlgC}{\text{\tbf{Alg}}_{\C}} % cat of complex algebras
\DeclareMathOperator*{\bAlgK}{\text{\tbf{Alg}}_{\K}} % cat of K-algebras

\DeclareMathOperator{\leqalg}{\text{Alg}} % cat of subalgebras 
% ######################################################
% ######################################################
% ######################################################





% Universal Algebra
\DeclareMathOperator*{\ar}{\text{arity}} % arity
\DeclareMathOperator*{\Sg}{\text{Sg}} 
\DeclareMathOperator{\uni}{Uni}
\DeclareMathOperator{\oper}{Oper}
\DeclareMathOperator{\indoper}{IndOper}
\DeclareMathOperator{\typ}{Type}
\DeclareMathOperator{\typs}{Types}
\DeclareMathOperator{\typfun}{TypeFun}
\DeclareMathOperator{\subu}{SubUni}
\DeclareMathOperator{\suba}{SubAlg}



























% ###################_________Order Theory_________###################


\DeclareMathOperator{\upar}{\uparrow}
\DeclareMathOperator{\doar}{\downarrow}
\DeclareMathOperator{\PosRefl}{PosRefl}

\DeclareMathOperator{\Mtn}{\text{Mtn}} % set of monotone maps
\DeclareMathOperator{\Antn}{\text{Antn}} % set of antitone maps

\DeclareMathOperator{\RAdj}{RAdj}
\DeclareMathOperator{\rad}{\text{rad}}
\DeclareMathOperator{\lad}{\text{lad}}
\DeclareMathOperator{\RQAdj}{RQAdj}
\DeclareMathOperator{\LAdj}{LAdj}
\DeclareMathOperator{\LQAdj}{LQAdj}

\DeclareMathOperator{\RConn}{RConn}
\DeclareMathOperator{\LConn}{LConn}

\DeclareMathOperator{\Ideals}{Ideals}
\DeclareMathOperator{\Filters}{Filters}
\DeclareMathOperator{\IdealGen}{IdealGen}
\DeclareMathOperator{\FilterGen}{FilterGen}


\DeclareMathOperator*{\prmj}{\text{prime}_{\vee}} % join prime
\DeclareMathOperator*{\prmm}{\text{prime}_{\wedge}} % meet prime
\DeclareMathOperator*{\irrj}{\text{irred}_{\vee}} % join irreducible
\DeclareMathOperator*{\irrm}{\text{irred}_{\wedge}} % meet irreducible

\DeclareMathOperator{\supSet}{supSet} % supremum set
\DeclareMathOperator{\infSet}{infSet} % infimum set

\DeclareMathOperator{\lbd}{lb}
\DeclareMathOperator{\ubd}{ub}


% Classes

% Categories
\DeclareMathOperator*{\bPos}{\text{\tbf{Pos}}} % cat of posets
























% ###################_________Category Theory________###################

\newcommand{\opp}{\text{op}} % opposite functor

\DeclareMathOperator{\id}{id} % identity
\DeclareMathOperator{\Hom}{Hom} % hom set
\DeclareMathOperator{\Endo}{End} % endomorphisms
\DeclareMathOperator{\Iso}{Iso} % ismorphisms
\DeclareMathOperator{\Aut}{Aut} % automorphisms
\DeclareMathOperator{\Mono}{Mono} % monomorphisms
\DeclareMathOperator{\ExtMono}{ExtMono} % extremal monomorphisms
\DeclareMathOperator{\RegMono}{RegMono} % regular monomorphisms
\DeclareMathOperator{\Epi}{Epi} % epimorphisms
\DeclareMathOperator{\ExtEpi}{ExtEpi} % extremal monomorphisms
\DeclareMathOperator{\RegEpi}{RegEpi} % regular monomorphisms

%\DeclareMathOperator{\Term}{Term} % terminal objects in a category
\DeclareMathOperator{\Coterm}{Coterm} % coterminal objects in a category

\DeclareMathOperator{\Prod}{Prod} % product objects in a category
\DeclareMathOperator{\Coprod}{Coprod} % coproduct object in a category

\DeclareMathOperator{\Exp}{Exp} % exponential objects in a category
\DeclareMathOperator{\Coexp}{Coexp} % coexponential objects in a category

\DeclareMathOperator{\Tmnl}{Term} % terminal objects in a category
\DeclareMathOperator{\Init}{Init} % initial objects in a category

\DeclareMathOperator{\Eqzr}{Equi} % equilizer objects in a category
\DeclareMathOperator{\Coeqzr}{Coequi} % coequilizer objects in a category



\DeclareMathOperator{\bMono}{\tbf{Mono}} % subcategory restricted to monomorphisms
\DeclareMathOperator{\bEpi}{\tbf{Epi}} % subcategory restricted to monomorphisms

\DeclareMathOperator{\Sub}{Sub} % subobject proset

\DeclareMathOperator{\colim}{colim} % colimit of diagram
\DeclareMathOperator{\limSet}{limSet} % limit set of diagram
\DeclareMathOperator{\colimSet}{colimSet} % colimit set of diagram
\DeclareMathOperator{\limChoice}{limChoice} % limit choice function 
\DeclareMathOperator{\colimChoice}{colimChoice} % colimit choice function 

% Cateogries
\DeclareMathOperator{\bArr}{\textbf{Arr}} % arrow category
\DeclareMathOperator*{\bCat}{\text{\tbf{Cat}}} % category of small categories
\DeclareMathOperator*{\bCone}{\text{\tbf{Cone}}} % category of small cones
\DeclareMathOperator*{\bCocone}{\text{\tbf{Cocone}}} % category of small cocones


































\newcommand{\bQuiv}{\text{\tbf{Quiv}}}


\newcommand{\Factors}{\text{Factors}}
\newcommand{\Factor}{\text{Factor}}

\newcommand{\LFactors}{\text{LFactors}}
\newcommand{\LFactor}{\text{LFactor}}


\newcommand{\RFactors}{\text{RFactors}}
\newcommand{\RFactor}{\text{RFactor}}


\DeclareMathOperator*{\bBanAlg}{\text{\tbf{BanAlg}}}



\DeclareMathOperator*{\bNorm}{\text{\tbf{Norm}}} 
\DeclareMathOperator*{\bBan}{\text{\tbf{Ban}}} 
\DeclareMathOperator*{\bHilb}{\text{\tbf{Hilb}}} 
\DeclareMathOperator*{\bProb}{\text{\tbf{Prob}}} 

\DeclareMathOperator*{\bRep}{\text{\tbf{Rep}}} 
\DeclareMathOperator*{\bURep}{\text{\tbf{URep}}}





\DeclareMathOperator*{\bTopField}{\text{\tbf{TopField}}} 
\DeclareMathOperator*{\bTopMon}{\text{\tbf{TopMon}}} 
\DeclareMathOperator*{\bTopGrp}{\text{\tbf{TopGrp}}}
\DeclareMathOperator*{\bTopVect}{\text{\tbf{TopVect}}} 










%\newcommand{\TopMeas}{\text{TopMeas}}
%\newcommand{\Msrpos}{\text{Msr}_{+}}
%\newcommand{\TopMsrpos}{\text{TopMsr}_{+}}
%\newcommand{\TopRadMsrpos}{\text{TopRadMsr}_{+}}
%\newcommand{\TopRadMsrone}{\text{TopRadMsr}_{1}}
%\newcommand{\MsrC}{\text{Msr}_{\C}} 
%\newcommand{\TopMsrC}{\text{TopMsr}_{\C}} 
%\newcommand{\TopRadMsrC}{\text{TopRadMsr}_{\C}} 
%\newcommand{\Maninf}{\text{Man}^{\infty}} 
%\newcommand{\ManBndinf}{\text{ManBnd}^{\infty}} 
%\newcommand{\Man0}{\text{Man}^{0}}
%\newcommand{\Buninf}{\text{Bun}^{\infty}} 
%\newcommand{\VecBuninf}{\text{VecBun}^{\infty}} 
%
%\newcommand{\Mon}{\text{Mon}} 
%\newcommand{\Grp}{\text{Grp}}
%\newcommand{\Semgrp}{\text{Semgrp}}
%\newcommand{\LieGrp}{\text{LieGrp}} 
%
%\newcommand{\Field}{\text{Field}} 
%
%\newcommand{\Vect}{\text{Vect}}
%\newcommand{\Mod}{\text{Mod}}
%\newcommand{\Alg}{\text{Alg}} 
%
%\newcommand{\Rep}{\text{Rep}} 
%\newcommand{\URep}{\text{URep}}
%
%\newcommand{\bNorm}{\text{Norm}} 
%\newcommand{\Ban}{\text{Ban}} 
%\newcommand{\Hilb}{\text{Hilb}} 
%\newcommand{\BanAlg}{\text{BanAlg}}
%
%\newcommand{\Prob}{\text{Prob}} 
%\newcommand{\PrinBuninf}{\text{PrinBun}^{\infty}}
%
%\newcommand{\TopField}{\text{TopField}} 
%\newcommand{\TopMon}{\text{TopMon}}
%\newcommand{\TopGrp}{\text{TopGrp}}
%\newcommand{\TopVect}{\text{TopVect}} 
%\newcommand{\TopEq}{\text{TopEq}}
%\newcommand{\Frm}{\text{Frm}}
%\newcommand{\Loc}{\text{Loc}}  
%
%\newcommand{\VectR}{\text{Vect}_{\R}}
%\newcommand{\VectC}{\text{Vect}_{\C}} 
%\newcommand{\VectK}{\text{Vect}_{\K}}
%
%\newcommand{\Cat}{\text{Cat}
%\newcommand{\Cone}{\text{Cone}
%\newcommand{\Cocone}{\text{Cocone}


\newcommand{\0}{\mbf{0}}
\newcommand{\1}{\mbf{1}}
\newcommand{\2}{\mbf{2}}







\DeclareMathOperator*{\Fac}{\text{Factor}}




% notation
\renewcommand{\r}{\rangle}
\renewcommand{\l}{\langle}
\renewcommand{\div}{\text{div}}
\renewcommand{\Re}{\text{Re} \,}
\renewcommand{\Im}{\text{Im} \,}
\newcommand{\Img}{\text{Img} \,}
\newcommand{\grad}{\text{grad}}





% cat theory notation
\newcommand{\pred}{\text{pred}}
\newcommand{\comp}{\text{comp}}


\newcommand{\Bool}{\text{Bool}}


\newcommand{\tbf}[1]{\textbf{#1}}
\newcommand{\tcb}[1]{\textcolor{blue}{#1}}
\newcommand{\tcr}[1]{\textcolor{red}{#1}}
\newcommand{\mbf}[1]{\mathbf{#1}}


\newcommand{\ol}[1]{\overline{#1}}
\newcommand{\ub}[1]{\underbar{#1}}
\newcommand{\tl}[1]{\tilde{#1}}
\newcommand{\p}{\partial}
\newcommand{\Tn}[1]{T^{r_{#1}}_{s_{#1}}(V)}
\newcommand{\Tnp}{T^{r_1 + r_2}_{s_1 + s_2}(V)}
\newcommand{\Perm}{\text{Perm}}
\newcommand{\wh}[1]{\widehat{#1}}
\newcommand{\wt}[1]{\widetilde{#1}}
\newcommand{\defeq}{\vcentcolon=}

\newcommand{\KosConn}{\text{KosConn}}

\newcommand{\tri}{\triangle}
\newcommand{\trl}{\triangleleft}
\newcommand{\trr}{\triangleright}
\newcommand{\loz}{\lozenge}
\newcommand{\dmd}{\diamond}
\newcommand{\alg}{\text{alg}}
\newcommand{\Triv}{\text{Triv}}
\newcommand{\Der}{\text{Der}}


\DeclareMathOperator{\conj}{conj}
\DeclareMathOperator{\conjdom}{\text{dom}_{\text{conj}}}
\DeclareMathOperator{\biconj}{\text{biconj}}
\DeclareMathOperator{\biconjdom}{\text{dom}_{\text{biconj}}}



\newcommand{\lcm}{\text{lcm}}
\newcommand{\Imax}{\MI_{\text{max}}}


\DeclareMathOperator{\Rl}{\text{Re}}
\DeclareMathOperator{\Imn}{\text{Imn}}

\newcommand{\veq}{\rotatebox{90}{$\,=$}}
\newcommand{\dtbeq}{\rotatebox{-45}{$\,=$}}
\newcommand{\dbteq}{\rotatebox{45}{$\,=$}}




\newcommand{\eval}{\text{eval}}
\newcommand{\coeval}{\text{coeval}}




% limits



\newcommand{\limfn}{\liminf \limits_{n \rightarrow \infty}}
\newcommand{\limpn}{\limsup \limits_{n \rightarrow \infty}}

\newcommand{\limn}{\lim \limits_{n \rightarrow \infty}}
\newcommand{\limx}[1]{\lim \limits_{x \rightarrow #1}}
\newcommand{\limt}[1]{\lim \limits_{t \rightarrow #1}}
\newcommand{\limh}[1]{\lim \limits_{h \rightarrow #1}}

\newcommand{\limneq}[2]{\lim\limits_{\substack{#1 \rightarrow #2 \\ #1 \neq #2}}}
\newcommand{\limneqx}[1]{\lim\limits_{\substack{x \rightarrow #1 \\ x \neq #1}}}
\newcommand{\limneqt}[1]{\lim\limits_{\substack{t \rightarrow #1 \\ x \neq #1}}}

\newcommand{\limfneq}[2]{\liminf\limits_{\substack{{#1} \rightarrow #2 \\ #1 \neq #2}}}
\newcommand{\limfneqx}[1]{\liminf\limits_{\substack{x \rightarrow #1 \\ x \neq #1}}}

\newcommand{\convt}[1]{\xrightarrow{\text{#1}}}
\newcommand{\conv}[1]{\xrightarrow{#1}} 
\newcommand{\seq}[2]{(#1_{#2})_{#2 \in \N}}

% intervals
\newcommand{\Rext}{[-\infty,\infty]}


\newcommand{\RG}{[0,\infty]}
\newcommand{\Rg}{[0,\infty)}
\newcommand{\Rgp}{(0,\infty)}
\newcommand{\Ru}{(\infty, \infty]}
\newcommand{\Rd}{[\infty, \infty)}
\newcommand{\ui}{[0,1]}

\newcommand{\RE}{\ol{\R}}
\newcommand{\REp}{\ol{\R}_+}

\newcommand{\Rp}{\R_+}







% abreviations 
\newcommand{\lsc}{lower semicontinuous}

% misc
\newcommand{\as}[1]{\overset{#1}{\sim}}
\newcommand{\astx}[1]{\overset{\text{#1}}{\sim}}
\newcommand{\io}{\text{ i.o.}}
%\newcommand{\ev}{\text{ ev.}}
\newcommand{\Ll}{L^1_{\text{loc}}(\R^n)}

\newcommand{\loc}{\text{loc}}
\newcommand{\BV}{\text{BV}}
\newcommand{\NBV}{\text{NBV}}
\newcommand{\TV}{\text{TV}}

\newcommand{\op}[1]{{#1}^{\text{op}}}
\newcommand{\lop}{\leq^*}